\documentclass[12pt,a4paper]{exam}
\usepackage[utf8]{inputenc}
\usepackage{amsmath,amssymb,amsfonts}
\usepackage{geometry}
\usepackage{enumitem}
\usepackage{graphicx}
\usepackage{hyperref}
\usepackage{xcolor}
\geometry{a4paper, top=2cm, bottom=2cm, left=2.5cm, right=2.5cm}
\hypersetup{colorlinks=true, linkcolor=blue, urlcolor=blue}
\renewcommand{\thequestion}{\textbf{\arabic{question}}}
\renewcommand{\questionlabel}{\thequestion.}
\pointname{ marks}
\pointsdroppedatright
\bracketedpoints
\setlength{\parindent}{0pt}
\setlength{\emergencystretch}{3em}
\allowdisplaybreaks
\newsavebox{\schmathbox}
\newcommand{\fitmath}[1]{%
  \sbox{\schmathbox}{$\displaystyle #1$}%
  \ifdim\wd\schmathbox>\linewidth
    \resizebox{\linewidth}{!}{\usebox{\schmathbox}}%
  \else
    \usebox{\schmathbox}%
  \fi}

\begin{document}

\thispagestyle{empty}
\begin{center}
  \textbf{\LARGE Diag} \\[0.3cm]
  \textbf{Time Allowed: 3 Hours} \hfill \textbf{Maximum Marks: 80} \\[0.3cm]
  \rule{\textwidth}{0.5pt}
\end{center}

\vspace{0.3cm}

\noindent\textbf{General Instructions:}
\begin{enumerate}
  \item {\normalfont\normalcolor This question paper contains 40 questions. All questions are compulsory.}
  \item {\normalfont\normalcolor Questions are grouped into sections A through E.}
  \item {\normalfont\normalcolor Use of calculators is NOT permitted.}
\end{enumerate}
\bigskip
\noindent\textbf{\large Section A: Multiple Choice \& Assertion-Reason (1 mark each)}

\begin{questions}
\question[1] If $p(x) = g(x) \cdot q(x) + r(x)$, and $p(x) = x^2 + 3x + 2$, $g(x) = x + 1$, then the quotient $q(x)$ is:
\begin{enumerate}[label=(\Alph*)]
  \item (A) x + 2
  \item (B) x - 2
  \item (C) x + 1
  \item (D) x - 1
\end{enumerate}
\vspace{0.15cm}

\question[1] According to the division algorithm for polynomials, if $p(x)$ is divided by $g(x)$, then the remainder $r(x)$ must satisfy:
\begin{enumerate}[label=(\Alph*)]
  \item (A) deg(r(x)) > deg(g(x))
  \item (B) deg(r(x)) = deg(g(x))
  \item (C) deg(r(x)) < deg(g(x))
  \item (D) r(x) = 1
\end{enumerate}
\vspace{0.15cm}

\question[1] If $p(x) = x^2 - 4$ is divided by $g(x) = x - 2$, what is the remainder?
\begin{enumerate}[label=(\Alph*)]
  \item (A) 4
  \item (B) 2
  \item (C) -4
  \item (D) 0
\end{enumerate}
\vspace{0.15cm}

\question[1] Given $p(x) = 2x^2 + 5x + 3$ and $g(x) = 2x + 3$, the quotient $q(x)$ is:
\begin{enumerate}[label=(\Alph*)]
  \item (A) x - 1
  \item (B) x + 1
  \item (C) 2x + 1
  \item (D) x + 3
\end{enumerate}
\vspace{0.15cm}

\question[1] If the dividend is $x^2 + 5x + 6$, the divisor is $x + 2$, and the remainder is $0$, what is the quotient?
\begin{enumerate}[label=(\Alph*)]
  \item (A) x + 2
  \item (B) x - 3
  \item (C) x + 3
  \item (D) x - 2
\end{enumerate}
\vspace{0.15cm}

\question[1] A tangent $PQ$ at a point $P$ of a circle of radius $5$ cm meets a line through the center $O$ at a point $Q$ so that $OQ = 13$ cm. The length of $PQ$ is:
\begin{enumerate}[label=(\Alph*)]
  \item (A) 10 cm
  \item (B) 11 cm
  \item (C) 12 cm
  \item (D) 15 cm
\end{enumerate}
\vspace{0.15cm}

\question[1] If tangents $PA$ and $PB$ from a point $P$ to a circle with center $O$ are inclined to each other at an angle of $80^\circ$, then $\angle POA$ is equal to:
\begin{enumerate}[label=(\Alph*)]
  \item $(A) 50^\circ$
  \item $(B) 60^\circ$
  \item $(C) 70^\circ$
  \item $(D) 80^\circ$
\end{enumerate}
\vspace{0.15cm}

\question[1] The number of tangents that can be drawn from a point lying inside a circle is:
\begin{enumerate}[label=(\Alph*)]
  \item (A) 0
  \item (B) 1
  \item (C) 2
  \item (D) Infinite
\end{enumerate}
\vspace{0.15cm}

\question[1] Assertion (A): For a polynomial $p(x) = x^2 - 3x + 2$ and divisor $g(x) = x - 1$, the remainder is 0. Reason (R): If $g(x)$ is a factor of $p(x)$, then $p(x) = g(x)q(x) + r(x)$, where $r(x) = 0$.
\begin{enumerate}[label=(\Alph*)]
  \item (A) Both Assertion (A) and Reason (R) are true and Reason (R) is the correct explanation of Assertion (A)
  \item (B) Both Assertion (A) and Reason (R) are true but Reason (R) is NOT the correct explanation of Assertion (A)
  \item (C) Assertion (A) is true but Reason (R) is false
  \item (D) Assertion (A) is false but Reason (R) is true
\end{enumerate}
\vspace{0.15cm}

\question[1] Assertion (A): The degree of the remainder $r(x)$ must be less than the degree of the divisor $g(x)$. Reason (R): In the division algorithm $p(x) = g(x)q(x) + r(x)$, $r(x)$ can be a non-zero constant even if $g(x)$ has degree 1.
\begin{enumerate}[label=(\Alph*)]
  \item (A) Both Assertion (A) and Reason (R) are true and Reason (R) is the correct explanation of Assertion (A)
  \item (B) Both Assertion (A) and Reason (R) are true but Reason (R) is NOT the correct explanation of Assertion (A)
  \item (C) Assertion (A) is true but Reason (R) is false
  \item (D) Assertion (A) is false but Reason (R) is true
\end{enumerate}
\vspace{0.15cm}

\question[1] Assertion (A): If $p(x) = x^2 + 1$ is divided by $g(x) = x$, the remainder is 1. Reason (R): The remainder of $p(x)$ divided by $x-a$ is $p(a)$.
\begin{enumerate}[label=(\Alph*)]
  \item (A) Both Assertion (A) and Reason (R) are true and Reason (R) is the correct explanation of Assertion (A)
  \item (B) Both Assertion (A) and Reason (R) are true but Reason (R) is NOT the correct explanation of Assertion (A)
  \item (C) Assertion (A) is true but Reason (R) is false
  \item (D) Assertion (A) is false but Reason (R) is true
\end{enumerate}
\vspace{0.15cm}

\question[1] Assertion (A): If $p(x) = 2x^2 + 4x + 2$ is divided by $g(x) = x+1$, the quotient is $2x+2$. Reason (R): $2x^2 + 4x + 2 = (x+1)(2x+2) + 0$.
\begin{enumerate}[label=(\Alph*)]
  \item (A) Both Assertion (A) and Reason (R) are true and Reason (R) is the correct explanation of Assertion (A)
  \item (B) Both Assertion (A) and Reason (R) are true but Reason (R) is NOT the correct explanation of Assertion (A)
  \item (C) Assertion (A) is true but Reason (R) is false
  \item (D) Assertion (A) is false but Reason (R) is true
\end{enumerate}
\vspace{0.15cm}

\question[1] Assertion (A): The division algorithm for polynomials is applicable only when the divisor $g(x) \neq 0$. Reason (R): Division by zero is undefined in mathematics.
\begin{enumerate}[label=(\Alph*)]
  \item (A) Both Assertion (A) and Reason (R) are true and Reason (R) is the correct explanation of Assertion (A)
  \item (B) Both Assertion (A) and Reason (R) are true but Reason (R) is NOT the correct explanation of Assertion (A)
  \item (C) Assertion (A) is true but Reason (R) is false
  \item (D) Assertion (A) is false but Reason (R) is true
\end{enumerate}
\vspace{0.15cm}

\end{questions}
\bigskip
\noindent\textbf{\large Section B: Very Short Answer (2 marks each)}

\begin{questions}
\question[2] On dividing $x^3 - 3x^2 + x + 2$ by a polynomial $g(x)$, the quotient and remainder were $x - 2$ and $-2x + 4$ respectively. Find $g(x)$.
\vspace{0.15cm}

\question[2] Check whether $x^2 + 3x + 1$ is a factor of $3x^4 + 5x^3 - 7x^2 + 2x + 2$ by applying the division algorithm.
\vspace{0.15cm}

\question[2] Find the quotient and remainder when $p(x) = x^3 - 3x^2 + 5x - 3$ is divided by $g(x) = x^2 - 2$.
\vspace{0.15cm}

\question[2] What must be subtracted from $f(x) = x^4 + 2x^3 - 2x^2 + 4x - 1$ so that the resulting polynomial is exactly divisible by $x^2 + 2$?
\vspace{0.15cm}

\question[2] If the polynomial $x^4 - 6x^3 + 16x^2 - 25x + 10$ is divided by another polynomial $x^2 - 2x + k$, the remainder comes out to be $x + a$. Find the values of $k$ and $a$.
\vspace{0.15cm}

\end{questions}
\bigskip
\noindent\textbf{\large Section C: Short Answer (3 marks each)}

\begin{questions}
\question[3] The mean of the following frequency distribution is $50$. Find the value of $p$ if the sum of frequencies is $100$: 
| Class | 0-20 | 20-40 | 40-60 | 60-80 | 80-100 |
|---|---|---|---|---|---|
| Frequency | 17 | $p$ | 32 | $p$ | 19 |
\vspace{0.15cm}

\question[3] The following distribution shows the daily pocket allowance of children of a locality. If the mean pocket allowance is Rs. $18$, find the missing frequency $f$:
| Daily allowance (in Rs) | 11-13 | 13-15 | 15-17 | 17-19 | 19-21 | 21-23 | 23-25 |
|---|---|---|---|---|---|---|---|
| Number of children | 7 | 6 | 9 | 13 | $f$ | 5 | 4 |
\vspace{0.15cm}

\question[3] Find the median of the following data:
| Marks | 0-10 | 10-20 | 20-30 | 30-40 | 40-50 |
|---|---|---|---|---|---|
| Number of students | 5 | 15 | 20 | 7 | 3 |
\vspace{0.15cm}

\question[3] The following table shows the ages of 100 patients admitted in a hospital during a year. Find the mode of the ages:
| Age (in years) | 5-15 | 15-25 | 25-35 | 35-45 | 45-55 | 55-65 |
|---|---|---|---|---|---|---|
| Number of patients | 6 | 11 | 21 | 23 | 14 | 5 |
\vspace{0.15cm}

\question[3] A survey conducted on 20 households in a locality by a group of students resulted in the following frequency table for the number of family members in a household:
| Family size | 1-3 | 3-5 | 5-7 | 7-9 | 9-11 |
|---|---|---|---|---|---|
| Frequency | 7 | 8 | 2 | 2 | 1 | 
Find the mode of this data.
\vspace{0.15cm}

\question[3] On dividing $x^3 - 3x^2 + x + 2$ by a polynomial $g(x)$, the quotient and remainder were $x - 2$ and $-2x + 4$, respectively. Find $g(x)$.
\vspace{0.15cm}

\question[3] Find the values of $a$ and $b$ if the polynomial $x^4 + x^3 + 8x^2 + ax + b$ is exactly divisible by $x^2 + 1$.
\vspace{0.15cm}

\question[3] Divide the polynomial $p(x) = 3x^2 - x^3 - 3x + 5$ by $g(x) = x - 1 - x^2$ and verify the division algorithm.
\vspace{0.15cm}

\question[3] What must be subtracted from $8x^4 + 14x^3 - 2x^2 + 7x - 8$ so that the resulting polynomial is exactly divisible by $4x^2 + x - 2$?
\vspace{0.15cm}

\question[3] Obtain all zeros of $3x^4 + 6x^3 - 2x^2 - 10x - 5$, if two of its zeros are $\sqrt{5/3}$ and $-\sqrt{5/3}$.
\vspace{0.15cm}

\end{questions}
\bigskip
\noindent\textbf{\large Section D: Long Answer (4-5 marks each)}

\begin{questions}
\question[5] The mean of the following frequency distribution is $50$. Find the values of $p$ and $q$, if the sum of all frequencies is $120$. 

| Class Interval | 0-20 | 20-40 | 40-60 | 60-80 | 80-100 |
| :--- | :--- | :--- | :--- | :--- | :--- |
| Frequency | 17 | $p$ | 32 | $q$ | 19 |
\vspace{0.15cm}

\question[5] The median of the following data is $525$. Find the values of $x$ and $y$, if the total frequency is $100$.

| Class Interval | 0-100 | 100-200 | 200-300 | 300-400 | 400-500 | 500-600 | 600-700 | 700-800 | 800-900 | 900-1000 |
| :--- | :--- | :--- | :--- | :--- | :--- | :--- | :--- | :--- | :--- | :--- |
| Frequency | 2 | 5 | $x$ | 12 | 17 | 20 | $y$ | 9 | 7 | 4 |
\vspace{0.15cm}

\question[5] The following distribution gives the daily income of $50$ workers of a factory. Convert the distribution above to a 'less than type' cumulative frequency distribution and draw its ogive. From the ogive, find the median daily income.

| Daily Income (in Rs.~) | 100-120 | 120-140 | 140-160 | 160-180 | 180-200 |
| :--- | :--- | :--- | :--- | :--- | :--- |
| Number of workers | 12 | 14 | 8 | 6 | 10 |
\vspace{0.15cm}

\question[5] Find the mode of the following frequency distribution. If the mean of the same data is $28.6$, verify the empirical relationship between mean, median, and mode.

| Class | 0-10 | 10-20 | 20-30 | 30-40 | 40-50 |
| :--- | :--- | :--- | :--- | :--- | :--- |
| Frequency | 5 | 8 | 20 | 15 | 7 |
\vspace{0.15cm}

\question[5] Obtain all zeroes of the polynomial $p(x) = 2x^4 - 3x^3 - 3x^2 + 6x - 2$, if two of its zeroes are $\sqrt{2}$ and $-\sqrt{2}$.
\vspace{0.15cm}

\question[5] Find the values of $a$ and $b$ so that the polynomial $x^4 + x^3 + 8x^2 + ax + b$ is exactly divisible by $x^2 + 1$.
\vspace{0.15cm}

\question[5] If the polynomial $f(x) = x^4 - 6x^3 + 16x^2 - 25x + 10$ is divided by another polynomial $g(x) = x^2 - 2x + k$, the remainder comes out to be $x + a$. Find the values of $k$ and $a$.
\vspace{0.15cm}

\question[5] A company models its profit $P(x)$ in thousands of rupees by the polynomial $P(x) = -x^3 + 4x^2 + x - 4$, where $x$ represents the number of units sold. If the factor representing the price per unit is $(x-4)$, find the other factors of the profit polynomial.
\vspace{0.15cm}

\question[5] Prove that the lengths of tangents drawn from an external point to a circle are equal.
\vspace{0.15cm}

\question[5] A quadrilateral $\text{ABCD}$ is drawn to circumscribe a circle. Prove that $AB + CD = AD + BC$.
\vspace{0.15cm}

\question[5] Two concentric circles are of radii 5 cm and 3 cm. Find the length of the chord of the larger circle which touches the smaller circle.
\vspace{0.15cm}

\question[5] If a hexagon $\text{ABCDEF}$ circumscribes a circle, prove that $AB + CD + EF = BC + DE + FA$.
\vspace{0.15cm}

\end{questions}
\bigskip

\end{document}
